\documentclass[11pt,a4paper]{article}
\usepackage[margin=0.9in]{geometry}
\usepackage{amsmath,amssymb}
\usepackage{booktabs}
\usepackage{threeparttable}
\usepackage{array}
\usepackage{adjustbox}
\usepackage{caption}
\usepackage{float}
\usepackage{microtype}
\usepackage{setspace}
\setstretch{1.05}
\captionsetup{font=small,labelfont=bf}
\begin{document}

\appendix
\section{Appendix: Additional Details for the Indirect-Inference Exercise}

This appendix records the details that should not be loaded into the main text. The main section should remain focused on the causal-core message: the model reproduces the output-competition IV response and passes a residual-IV validation. The appendix should instead document the estimator, the optimization path, non-targeted diagnostics, and the reasons why the current run should not yet be interpreted as a full welfare model.

\subsection{Estimator and auxiliary moments}

The estimation is an indirect-inference exercise because the model is disciplined by auxiliary empirical objects rather than by a closed-form likelihood. For each candidate parameter vector \(\theta\), the model generates a simulated panel with the same observation structure as the data. The same auxiliary procedures are then applied to the observed and simulated panels. In the notation of indirect inference, the binding function is
\begin{equation}
    h(\theta)=\operatorname{plim} m^{s}(\theta),
\end{equation}
where \(m^{s}(\theta)\) is the vector of simulated auxiliary moments. Identification requires that the chosen moments move differently when the economically relevant components of \(\theta\) are perturbed. In the present run, identification comes from six output-competition moments:
\begin{enumerate}
    \item the baseline IV response of incumbent markup growth;
    \item grouped-IV effects at q80, q85, and q90;
    \item the q90 interaction with sectoral concentration; and
    \item the sector inverse-markup IV response.
\end{enumerate}
The criterion is
\begin{equation}
Q(\theta)=\left[\widehat m-m^{s}(\theta)\right]'W
\left[\widehat m-m^{s}(\theta)\right],
\end{equation}
with \(W\) chosen to emphasize statistically stable causal moments. The empirical interpretation of the fitted model therefore depends on the content of these auxiliary moments. It is a model of the output-competition mechanism, not a model of every observed movement in markups.

\subsection{Full parameter list}

\begin{table}[H]
\centering
\caption{Full parameter vector in the causal-core run}
\begin{threeparttable}
\small
\begin{tabular}{lrlp{7.2cm}}
\toprule
Parameter & Value & Status & Interpretation \\
\midrule
\texttt{output\_scale} & 3.449 & estimated & Scale of the output-competition channel \\
\texttt{output\_markup} & -0.674 & estimated & Baseline output-shock effect on markup growth \\
\texttt{high\_markup\_output} & -1.514 & estimated & Extra output-shock effect for high-markup firms or sectors \\
\texttt{concentration\_output} & -3.964 & estimated & Extra output-shock effect in concentrated markets \\
\texttt{share\_output} & 0.500 & estimated & Output-shock component in market-share dynamics \\
\texttt{share\_input} & 0.000 & fixed & Input channel switched off \\
\texttt{input\_markup} & 0.000 & fixed & Input-markup channel switched off \\
\texttt{mean\_reversion} & 1.382 & estimated & Pullback of markups toward reference level \\
\texttt{drift} & -0.104 & estimated & Average residual drift \\
\texttt{exit\_intercept} & -2.300 & fixed & Baseline exit block fixed \\
\texttt{exit\_ip} & 0.000 & fixed & No output-shock selection effect \\
\texttt{exit\_markup} & 0.000 & fixed & No markup-dependent exit effect \\
\texttt{eta} & 1.010 & fixed & Elasticity parameter \\
\texttt{nu} & 4.000 & fixed & Elasticity parameter \\
\texttt{rho} & 8.000 & fixed & Elasticity parameter \\
\texttt{lambda\_min} & 0.050 & fixed & Lower bound for lambda \\
\texttt{lambda\_max} & 0.980 & fixed & Upper bound for lambda \\
\texttt{default\_lambda} & 0.850 & fixed & Default lambda \\
\bottomrule
\end{tabular}
\end{threeparttable}
\end{table}

The fixed zero input and selection parameters are important for interpretation. They explain why the model should be expected to miss input-shock and exit diagnostics. Those misses should guide the next structural extension, not overturn the targeted output-competition result.

\subsection{Optimization and targeted fit}

The optimizer made 512 objective evaluations. The objective declined from approximately \(37.635\) to \(1.727\). This is a substantial improvement and delivers economically interpretable signs: output competition lowers markup growth, and the effect is stronger where market power is initially high. The final objective is not close to zero because the current parameterization cannot perfectly trace the upper-tail gradient and sector aggregation response. This is acceptable for a causal-core run, provided the limitation is made explicit.

\begin{table}[H]
\centering
\caption{Moment gaps entering the objective}
\small
\begin{tabular}{lrrrr}
\toprule
Target moment & Data & Model & Gap & Weight \\
\midrule
\texttt{baseline\_iv\_dln\_mu} & -0.429 & -0.441 & 0.012 & 50.832 \\
\texttt{grouped\_iv\_q80} & -0.803 & -0.803 & -0.001 & 31.180 \\
\texttt{grouped\_iv\_q85} & -1.236 & -0.942 & -0.294 & 11.768 \\
\texttt{grouped\_iv\_q90} & -1.625 & -1.771 & 0.146 & 7.608 \\
\texttt{grouped\_iv\_q90\_cr4\_interaction} & -5.371 & -5.042 & -0.329 & 0.369 \\
\texttt{sector\_inverse\_markup\_iv} & 0.128 & 0.238 & -0.110 & 40.989 \\
\bottomrule
\end{tabular}
\end{table}

The q85 and q90 pattern is useful diagnostically. The model can generate upper-tail discipline, but the current interaction structure is too rigid to match the exact slope across the upper tail. A future version could allow smoother state dependence in initial markups and concentration rather than relying on a small number of coarse interaction terms.

\subsection{Residual-IV validation}

The residual-IV validation should be reported near the main SMM results rather than in a generic robustness section. Let
\begin{equation}
    r_{ijt}=\Delta\ln\mu_{ijt}-\Delta\ln\mu^{model}_{ijt}.
\end{equation}
define residual markup growth after removing the fitted output-competition component. Re-estimating the IV regression with \(r_{ijt}\) as the dependent variable gives the following validation result.

\begin{table}[H]
\centering
\caption{Residual-IV validation}
\small
\begin{tabular}{lrrrrr}
\toprule
Validation regression & Coef. & SE & Obs. & Clusters & First-stage \(F\) \\
\midrule
Residual markup growth on output shock & 0.042 & 0.250 & 31,375 & 90 & 22.930 \\
\bottomrule
\end{tabular}
\end{table}

The coefficient is small relative to its standard error and the first stage remains strong. This supports the interpretation that the model absorbs the output-shock variation it was designed to explain. The validation should not be overstated: it does not validate missing input, selection, or decomposition blocks.

\subsection{Non-targeted diagnostics}

The following diagnostics should be placed in the appendix. They are informative because they show where the model fails once it is asked to explain mechanisms outside the targeted output-competition channel.

\begin{table}[H]
\centering
\caption{Selected non-targeted diagnostic moments}
\small

\begin{adjustbox}{max width=\textwidth}
\begin{tabular}{lrrrl}
\toprule
Diagnostic moment & Data & Model & Gap & Interpretation \\
\midrule
\texttt{input\_shock\_between\_reallocation} & 8.055 & 3.744 & 4.310 & Underpredicts input-shock reallocation \\
\texttt{input\_shock\_within\_dln\_mu} & 0.142 & -0.306 & 0.448 & Wrong sign for input-markup relation \\
\texttt{exit\_rate} & 0.069 & 0.091 & -0.022 & Model exit rate too high \\
\texttt{exit\_iv} & 0.960 & 0.324 & 0.636 & Underpredicts output-shock selection response \\
\bottomrule
\end{tabular}
\end{adjustbox}
\end{table}

These misses are not surprising, because the input and selection channels are switched off. They are nevertheless economically important. The data contain an input-shock reallocation pattern and an exit response that the current model cannot reproduce. Therefore, the model should not be used yet for welfare analysis, long-run industry composition, or full manufacturing-wide counterfactuals.

\subsection{Counterfactual interpretation}

The no-output counterfactual removes the fitted output-competition channel. At the aggregate manufacturing level, the differences are small: the average inverse markup is \(0.8361\) in the observed model path and \(0.8350\) in the no-output counterfactual; the average aggregate markup is \(1.1966\) versus \(1.1981\). The 2011--2019 change in inverse markups is also close: \(-0.0183\) in the observed path and \(-0.0192\) in the no-output counterfactual.

The right interpretation is therefore local and sectoral. The model says that output competition disciplines exposed incumbent markups, especially in high-rent and concentrated markets. It does not say that Chinese competition explains most aggregate markup dynamics in Turkish manufacturing. This distinction is a strength of the exercise because it prevents the model from being used beyond the variation that identifies it.

\subsection{Recommended reporting allocation}

The main text should report only the objects needed for the causal-core claim:
\begin{enumerate}
    \item the motivation for using a model after IV;
    \item the SMM/indirect-inference criterion;
    \item the compact parameter table;
    \item the targeted causal moment fit;
    \item the residual-IV validation; and
    \item a restrained counterfactual synthesis.
\end{enumerate}
The appendix should report the full parameter vector, optimization history, full moment gaps, non-targeted diagnostics, input and selection failures, sector-level counterfactual heterogeneity, and allocative-wedge checks. The final thesis claim should remain: the exercise supports an output-competition markup-discipline channel, not a complete welfare model of Turkish manufacturing.

\end{document}
